
\subsection{Thermodynamic consistency and self-heating under cyclic loading: internal field analysis}

Upon the validated framework and mesh convergence in the previous benchmark, this second example investigates the transient thermodynamic response of the prestretched membrane under cyclic actuation. While standard experimental setups are typically restricted to macroscopic kinematic and electrical measurements, the proposed numerical framework provides high-fidelity access to internal thermodynamic fields that are difficult or impossible to observe in conventional laboratory settings.

The primary technical objectives of this second numerical example are:

\begin{enumerate}
    \item To assess the thermodynamic consistency of the coupled thermo-mechanical scheme by analyzing the evolution and conservation of internal energy, dissipation, and entropy over multiple loading cycles.
    \item To quantify the time-space variation of the temperature field induced by intrinsic dissipation and thermo-mechanical coupling (self-heating) under cyclic actuation.
    \item To demonstrate the predictive capability of the formulation in capturing complex, path-dependent internal variables that extend beyond standard experimentally observable macroscopic metrics.
\end{enumerate}


\begin{figure}[htpb]
  \centering
  \includegraphics[width=.95\textwidth]{cyclic/time_temperature.png}
  \caption{Evolution of the internal temperature field driven by thermo-mechanical self-heating over consecutive loading cycles.}
  \label{fig:time_temperature}
\end{figure}

